\section{Conclusion}
\label{sec:conclusion}

The boundary-layer lemma turns two enormous combined subgroup-transition
questions into 944,200 concrete physical obligations.  Exact rewriting and
neural-guided beam cascades then provide a bounded word for every obligation,
while deterministic
coverage and replay keep all learned components outside the proof base.  The
result is the unconditional pair of ambient-FTM upper bounds
$\Rsolve_{\Omega}(G_5,G_7)\leq21$ and
$\Rsolve_{\Omega}(G_7,G_9)\leq25$, and the conditional
global arithmetic 112---or 111 when combined with the separate color-neutral
first-phase reduction.

The broader lesson is methodological.  Learned search can participate in an
exhaustive computational proof without itself being complete, interpretable,
or numerically trusted, provided that its output is a collection of small
classical certificates, a deterministic checker proves exact coverage, and
deterministic quotient and physical-state replay validates every direct word.
The full-reproduction workflow additionally requires a separately implemented
Go replay of the same transcript.  In the published run, the survivor cascade concentrated the widest
searches on the final 29 roots of a $935{,}536$-root direct obligation set.
Determining its advantage over fixed-width or non-neural alternatives requires
a controlled comparison.  A natural next step is to study when
the same boundary-layer construction extends to thicker bands and whether
controlled cascade policies can predict the appropriate search budget before
running the full certificate pass.
